% English-language thesis example -- compile:  latexmk -xelatex example-en
\documentclass[english]{kuthesis}
% options: thai | english, numbered (เลขกำกับหัวข้อข้าง), guides (เส้นวัดระยะ),
%          latin=termes|inter|sarabun, math=termes|fira|lete|newcm|stix|none, thaifont=...

\usepackage[style=kuthesis]{biblatex}
\addbibresource{references.bib}

\thesistitle{การบำบัดน้ำเสียด้วยจุลินทรีย์แบบยึดเกาะตัวกลางชนิดเส้นใยไนล่อนโดยระบบถังกรองไร้อากาศ}
            {Wastewater Treatment of Fixed Film Microorganism on Nylon Rope Media by Anaerobic Filter}
\thaiauthor[นางสาว]{ซันวานี}{จิใจ}
\englishauthor[Miss]{Sunwanee}{Jijai}
\degree{วิศวกรรมศาสตรมหาบัณฑิต (วิศวกรรมสิ่งแวดล้อม)}
       {Master of Engineering (Environmental Engineering)}
\majorfield{วิศวกรรมสิ่งแวดล้อม}{Environmental Engineering}
\department{วิศวกรรมสิ่งแวดล้อม}{Environmental Engineering}
\advisor{รองศาสตราจารย์สมชาย ใจดี, Ph.D.}{Associate Professor Somchai Jaidee, Ph.D.}
\coadvisor{ผู้ช่วยศาสตราจารย์สมหญิง รักเรียน, D.Eng.}{Assistant Professor Somying Rakrian, D.Eng.}
\departmenthead{รองศาสตราจารย์สมศักดิ์ ตั้งใจ, Ph.D.}{Associate Professor Somsak Tangjai, Ph.D.}
\dean{รองศาสตราจารย์กัญจนา ธีระกุล, D.Agr.}{Associate Professor Gunjana Theeragool, D.Agr.}
\thesisyear{2549}{2006}
\ackdate{สิงหาคม 2549}{August 2006}
\kulogo[logo/ku-logo-en.png]{logo/ku-logo-th.png}  % ตรามหาวิทยาลัย (ไม่มีไฟล์ = วงกลมแทน)
\kuapprovallogo{logo/ku-logo-th-green.png}

\begin{document}

\frontmatter
\makecover
\makeapproval
\maketitlepage

\begin{thaiabstract}
  งานวิจัยนี้ศึกษาการบำบัดน้ำเสียด้วยถังกรองไร้อากาศ (anaerobic filter) ที่ใช้เส้นใยไนล่อนเป็นตัวกลางให้จุลินทรีย์ยึดเกาะ
  โดยเดินระบบที่ภาระบรรทุกสารอินทรีย์ 0.5--4.0 kg COD/m\textsuperscript{3}$\cdot$d
  ผลการทดลองพบว่าระบบสามารถกำจัดซีโอดีได้ร้อยละ 85--92

  ระยะเวลากักเก็บน้ำที่เหมาะสมคือ 12 ชั่วโมง และอัตราการเกิดก๊าซชีวภาพแปรผันตามภาระบรรทุกสารอินทรีย์
\end{thaiabstract}

\begin{englishabstract}
  This research studied an anaerobic filter using nylon rope as the support medium for
  attached microorganisms. The reactor was operated at organic loading rates of
  0.5--4.0~kg~COD/m\textsuperscript{3}$\cdot$d and removed 85--92\% of COD.

  The optimum hydraulic retention time was 12 hours.
\end{englishabstract}

\begin{acknowledgments}
  I would like to thank Associate Professor Dr. Somchai Jaidee, my thesis advisor, for his
  advice and encouragement throughout this work.

  I am grateful to my family for their support.
\end{acknowledgments}

\tableofcontents
\listoftables
\listoffigures

\begin{abbreviations}
  \abbr{BOD}{biochemical oxygen demand}
  \abbr{COD}{chemical oxygen demand}
  \abbr{HRT}{hydraulic retention time}
  \abbr{OLR}{organic loading rate}
  \abbr{$\mu$}{specific growth rate (d\textsuperscript{-1})}
\end{abbreviations}

\mainmatter

\chapter{Introduction}

Recently, industrial activities have experienced rapid growth producing an increase in the
volume and toxicity of residues \parencite{peterson2004}. \textcite{boddington2000a} reported
that agricultural practices affect mycorrhizal fungi, and similar results were found by
\textcite{chutima2551} and \textcite[3]{bago2000}. Several studies support this
\parencite{buddhakulsomsiri2008,seranevijaikitkhan2008,who2006}.

Biosorption, adsorption of heavy metal by biomass, has emerged as a potential alternative
to conventional techniques.

\chapter{Objectives}

\begin{enumerate}
  \item To study the COD removal efficiency of an anaerobic filter.
  \item To determine the optimum hydraulic retention time.
\end{enumerate}

\chapter{Literature Review}

\section{Anaerobic Treatment}

\subsection{Principles}

The specific growth rate follows the Monod equation (\ref{eq:monod}):
\begin{equation}
  \mu = \mu_{\max}\,\frac{S}{K_s + S} \label{eq:monod}
\end{equation}
where $S$ is the substrate concentration and $K_s$ the half-saturation constant.

\subsubsection{Temperature}
The optimum temperature for methanogens is 30--38\,°C.

\subsubsection{pH}
The optimum pH is 6.8--7.2.

\subsection{Support Media}

Media properties are given in Table~\ref{tab:media} and Figure~\ref{fig:reactor}.

\begin{table}[htbp]
  \caption{Properties of support media used in anaerobic filters.}
  \label{tab:media}
  \centering
  \begin{tabular}{lccc}
    \toprule
    Medium & Specific area & Porosity & Cost \\
           & (m\textsuperscript{2}/m\textsuperscript{3}) & (\%) & (Baht/m\textsuperscript{3}) \\
    \midrule
    Nylon rope   & 320 & 96 & 1,200 \\
    Plastic ring & 210 & 92 & 3,500 \\
    Gravel       & 60  & 45 & 400 \\
    \bottomrule
  \end{tabular}
  \tablesource{\textcite{peterson2004}}
\end{table}

\begin{figure}[htbp]
  \centering
  \fbox{\parbox[c][1.6in][c]{3.2in}{\centering (reactor diagram)}}
  \caption{Schematic diagram of the anaerobic filter with nylon rope media; a caption longer
           than one line is aligned with the first letter of the caption.}
  \label{fig:reactor}
\end{figure}

\chapter{Materials and Methods}
\section{Materials}
The reactor was made of clear acrylic.
\section{Methods}
The system was operated continuously for 180 days.

\chapter{Results and Discussion}
\section{Results}
Steady state was reached within 45 days.
\section{Discussion}
The results agree with \textcite{peterson2004}.

\chapter{Conclusion and Recommendations}
\section{Conclusion}
The nylon-rope anaerobic filter is efficient and inexpensive.
\section{Recommendations}
Operation at low temperature should be studied.

\printbibliography

\appendix
\appendixsection{Water Quality Analysis}

\begin{table}[htbp]
  \caption{Influent and effluent COD.}
  \centering
  \begin{tabular}{ccc}
    \toprule
    Day & Influent (mg/L) & Effluent (mg/L) \\
    \midrule
    1 & 2,450 & 380 \\
    2 & 2,510 & 342 \\
    \bottomrule
  \end{tabular}
\end{table}

\appendixsection{Data Sheets}

Data sheets used in the research.

\begin{vitae}
  \vitaeitem{NAME}{Miss Sunwanee Jijai}
  \vitaeitem{BIRTH DATE}{February 14, 1977}
  \vitaeitem{BIRTH PLACE}{Bangkok, Thailand}
  \vitaeitem{EDUCATION}{B.Eng. (Environmental Engineering), Kasetsart University, 1999}
  \vitaeitem{POSITION/TITLE}{Engineer}
  \vitaeitem{WORK PLACE}{Pollution Control Department}
  \vitaeitem{SCHOLARSHIP/AWARDS}{Graduate School teaching assistantship (2004)}
\end{vitae}

\end{document}
